\subsection{Decorators}

A decorator is a function transformation mechanism. It receives a function object, builds or selects another callable object, and returns that object to be bound under the original function name.

The important point is that decoration happens when the \texttt{def} statement executes, not when the decorated function is later called.

A decorator is not magic syntax around a function call. It is ordinary function-object passing, returning, and rebinding.

\subsubsection{Decorators as Function Transformation at Definition Time}

Start without decorator syntax.

\begin{verbatim}
def shout(text):
    return text.upper()

def plain(text):
    return text

print(f"shout:       {shout}")
print(f"type(shout): {type(shout)}")
print(f"shout('ni!'): {shout('ni!')}")
\end{verbatim}

Possible output:

\begin{verbatim}
shout:       <function shout at 0x...>
type(shout): <class 'function'>
shout('ni!'): NI!
\end{verbatim}

Now define a function that receives another function and returns a replacement function.

\begin{verbatim}
def make_logged(fn):
    def wrapper(text):
        print(f"calling {fn.__name__} with {text!r}")
        result = fn(text)
        print(f"result: {result!r}")
        return result

    return wrapper

def shout(text):
    return text.upper()

shout = make_logged(shout)

print(shout("no soup for you!"))
\end{verbatim}

Possible output:

\begin{verbatim}
calling shout with 'no soup for you!'
result: 'NO SOUP FOR YOU!'
NO SOUP FOR YOU!
\end{verbatim}

The critical line is:

\begin{verbatim}
shout = make_logged(shout)
\end{verbatim}

This line:

\begin{enumerate}
    \item reads the current function object bound to \texttt{shout};
    \item passes that function object into \texttt{make\_logged};
    \item receives the returned \texttt{wrapper} function object;
    \item rebinds the name \texttt{shout} to the wrapper.
\end{enumerate}

After that assignment, the name \texttt{shout} no longer points directly to the original function object. It points to the wrapper function object.

We can inspect that.

\begin{verbatim}
def make_logged(fn):
    def wrapper(text):
        print(f"calling {fn.__name__}")
        return fn(text)

    return wrapper

def shout(text):
    return text.upper()

original_shout = shout
shout = make_logged(shout)

print(f"original_shout: {original_shout}")
print(f"shout:          {shout}")
print(f"same object:    {original_shout is shout}")
print()
print(f"original name:  {original_shout.__name__}")
print(f"new name:       {shout.__name__}")
\end{verbatim}

Possible output:

\begin{verbatim}
original_shout: <function shout at 0x...>
shout:          <function make_logged.<locals>.wrapper at 0x...>
same object:    False

original name:  shout
new name:       wrapper
\end{verbatim}

The original function still exists because \texttt{original\_shout} and the wrapper closure still reference it. But the public name \texttt{shout} now refers to the wrapper.

The wrapper is a function object.

\begin{verbatim}
def make_logged(fn):
    def wrapper(text):
        print(f"calling {fn.__name__}")
        return fn(text)

    return wrapper

def shout(text):
    return text.upper()

shout = make_logged(shout)

print(f"type(shout): {type(shout)}")
print()

selected_names = [
    "__call__",
    "__name__",
    "__doc__",
    "__defaults__",
    "__code__",
    "__closure__",
]

for name in selected_names:
    print(f"{name:<16} {name in dir(shout)}")

print()
print(f"shout.__name__:    {shout.__name__!r}")
print(f"shout.__closure__: {shout.__closure__}")
\end{verbatim}

Possible output:

\begin{verbatim}
type(shout): <class 'function'>

__call__        True
__name__        True
__doc__         True
__defaults__    True
__code__        True
__closure__     True

shout.__name__:    'wrapper'
shout.__closure__: (<cell at 0x...: function object at 0x...>,)
\end{verbatim}

The wrapper has a closure because it remembers the original function object through the name \texttt{fn}.

The decorator function itself is also a normal function object.

\begin{verbatim}
def make_logged(fn):
    def wrapper(text):
        return fn(text)

    return wrapper

print(f"make_logged:       {make_logged}")
print(f"type(make_logged): {type(make_logged)}")
print(f"make_logged name:  {make_logged.__name__}")
print(f"make_logged code:  {make_logged.__code__}")
\end{verbatim}

Possible output:

\begin{verbatim}
make_logged:       <function make_logged at 0x...>
type(make_logged): <class 'function'>
make_logged name:  make_logged
make_logged code:  <code object make_logged at 0x..., file "...", line ...>
\end{verbatim}

A decorator can add behavior before the original call.

\begin{verbatim}
def announce(fn):
    def wrapper(text):
        print("Prepare yourself.")
        return fn(text)

    return wrapper

def say_quote(text):
    return text

say_quote = announce(say_quote)

print(say_quote("Nobody expects the Spanish Inquisition!"))
\end{verbatim}

Possible output:

\begin{verbatim}
Prepare yourself.
Nobody expects the Spanish Inquisition!
\end{verbatim}

A decorator can add behavior after the original call.

\begin{verbatim}
def add_exclamation(fn):
    def wrapper(text):
        result = fn(text)
        return result + "!"

    return wrapper

def say_quote(text):
    return text

say_quote = add_exclamation(say_quote)

print(say_quote("Ni"))
\end{verbatim}

Possible output:

\begin{verbatim}
Ni!
\end{verbatim}

A decorator can also block or replace the original call.

\begin{verbatim}
def block_call(fn):
    def wrapper(text):
        return "This function has been blocked."

    return wrapper

def say_quote(text):
    return text

say_quote = block_call(say_quote)

print(say_quote("D'oh!"))
\end{verbatim}

Possible output:

\begin{verbatim}
This function has been blocked.
\end{verbatim}

The wrapper does not have to call the original function. A decorator receives the original function object, but the returned callable decides what to do with it.

The structural pattern is:

\begin{verbatim}
def decorator(original_function):
    def replacement_function(...):
        ...
        result = original_function(...)
        ...
        return result

    return replacement_function

target_function = decorator(target_function)
\end{verbatim}

That is the basic decorator machine.

\subsubsection{The \texttt{@decorator} Syntax as Rebinding Sugar}

Python provides decorator syntax to write this transformation directly above the function definition.

This:

\begin{verbatim}
def make_logged(fn):
    def wrapper(text):
        print(f"calling {fn.__name__} with {text!r}")
        result = fn(text)
        print(f"result: {result!r}")
        return result

    return wrapper

@make_logged
def shout(text):
    return text.upper()

print(shout("no soup for you!"))
\end{verbatim}

Possible output:

\begin{verbatim}
calling shout with 'no soup for you!'
result: 'NO SOUP FOR YOU!'
NO SOUP FOR YOU!
\end{verbatim}

is equivalent to this:

\begin{verbatim}
def make_logged(fn):
    def wrapper(text):
        print(f"calling {fn.__name__} with {text!r}")
        result = fn(text)
        print(f"result: {result!r}")
        return result

    return wrapper

def shout(text):
    return text.upper()

shout = make_logged(shout)

print(shout("no soup for you!"))
\end{verbatim}

The decorator syntax is rebinding sugar.

The execution order is:

\begin{verbatim}
1. create the original function object from def shout(...)
2. pass that function object to make_logged
3. receive the returned wrapper object
4. bind the name shout to the wrapper object
\end{verbatim}

This happens when the \texttt{def} statement executes.

A print trace makes definition-time decoration visible.

\begin{verbatim}
def make_logged(fn):
    print(f"decorating {fn.__name__}")

    def wrapper(text):
        print(f"calling {fn.__name__}")
        return fn(text)

    return wrapper

print("before decorated def")

@make_logged
def shout(text):
    print("inside original shout")
    return text.upper()

print("after decorated def")
print()

print(shout("d'oh!"))
\end{verbatim}

Possible output:

\begin{verbatim}
before decorated def
decorating shout
after decorated def

calling shout
inside original shout
D'OH!
\end{verbatim}

The line:

\begin{verbatim}
decorating shout
\end{verbatim}

appears before the function is called. It appears while the decorated \texttt{def} statement is being executed.

This proves that decorator application is a definition-time event. The wrapper's body runs later, when the decorated function name is called.

Multiple decorators stack from bottom to top.

\begin{verbatim}
def deco_a(fn):
    print("applying deco_a")

    def wrapper(text):
        print("inside deco_a wrapper")
        return fn(text)

    return wrapper

def deco_b(fn):
    print("applying deco_b")

    def wrapper(text):
        print("inside deco_b wrapper")
        return fn(text)

    return wrapper

@deco_a
@deco_b
def say(text):
    print("inside original say")
    return text

print()
print(say("It is only a model."))
\end{verbatim}

Possible output:

\begin{verbatim}
applying deco_b
applying deco_a

inside deco_a wrapper
inside deco_b wrapper
inside original say
It is only a model.
\end{verbatim}

This:

\begin{verbatim}
@deco_a
@deco_b
def say(text):
    return text
\end{verbatim}

is equivalent to:

\begin{verbatim}
def say(text):
    return text

say = deco_a(deco_b(say))
\end{verbatim}

The decorator closest to the function is applied first. The outer decorator receives the result of the inner decoration.

Decorator expressions can also be calls. This creates a decorator factory.

\begin{verbatim}
def make_prefix_decorator(prefix):
    print(f"creating decorator for prefix={prefix!r}")

    def decorator(fn):
        print(f"decorating {fn.__name__}")

        def wrapper(text):
            result = fn(text)
            return f"{prefix}: {result}"

        return wrapper

    return decorator

@make_prefix_decorator("INFO")
def say(text):
    return text

print()
print(say("The answer is 42."))
\end{verbatim}

Possible output:

\begin{verbatim}
creating decorator for prefix='INFO'
decorating say

INFO: The answer is 42.
\end{verbatim}

This is equivalent to:

\begin{verbatim}
def say(text):
    return text

say = make_prefix_decorator("INFO")(say)
\end{verbatim}

There are two stages:

\begin{enumerate}
    \item \texttt{make\_prefix\_decorator("INFO")} returns a decorator;
    \item that returned decorator receives \texttt{say} and returns a wrapper.
\end{enumerate}

The returned decorator and wrapper are both function objects.

\begin{verbatim}
def make_prefix_decorator(prefix):
    def decorator(fn):
        def wrapper(text):
            return f"{prefix}: {fn(text)}"

        return wrapper

    return decorator

decorator_obj = make_prefix_decorator("INFO")

print(f"decorator_obj:       {decorator_obj}")
print(f"type(decorator_obj): {type(decorator_obj)}")
print(f"name:                {decorator_obj.__name__}")
print(f"closure:             {decorator_obj.__closure__}")
\end{verbatim}

Possible output:

\begin{verbatim}
decorator_obj:       <function make_prefix_decorator.<locals>.decorator at 0x...>
type(decorator_obj): <class 'function'>
name:                decorator
closure:             (<cell at 0x...: str object at 0x...>,)
\end{verbatim}

The decorator object remembers \texttt{prefix}. The wrapper object later remembers both \texttt{prefix} and \texttt{fn}.

The key rule is:

\begin{quote}
The \texttt{@decorator} form is not a new kind of function. It is syntax for passing the just-created function object through another callable and rebinding the original name to the returned object.
\end{quote}

\subsubsection{Wrapper Functions, Closure State, and Metadata Preservation}

A wrapper function often closes over the original function object.

\begin{verbatim}
def make_logged(fn):
    def wrapper(text):
        print(f"calling {fn.__name__}")
        return fn(text)

    return wrapper

@make_logged
def shout(text):
    """Return uppercase text."""
    return text.upper()

print(shout("ni!"))
print()
print(f"shout:          {shout}")
print(f"shout.__name__: {shout.__name__!r}")
print(f"shout.__doc__:  {shout.__doc__!r}")
print(f"closure:        {shout.__closure__}")
\end{verbatim}

Possible output:

\begin{verbatim}
calling shout
NI!

shout:          <function make_logged.<locals>.wrapper at 0x...>
shout.__name__: 'wrapper'
shout.__doc__:  None
closure:        (<cell at 0x...: function object at 0x...>,)
\end{verbatim}

The decorated name \texttt{shout} now points to \texttt{wrapper}. That is why \texttt{shout.\_\_name\_\_} is \texttt{'wrapper'} and the original docstring is not visible through \texttt{shout.\_\_doc\_\_}.

The original function object is preserved inside the wrapper's closure.

\begin{verbatim}
def make_logged(fn):
    def wrapper(text):
        print(f"calling {fn.__name__}")
        return fn(text)

    return wrapper

@make_logged
def shout(text):
    """Return uppercase text."""
    return text.upper()

cell = shout.__closure__[0]
original_fn = cell.cell_contents

print(f"cell:             {cell}")
print(f"type(cell):       {type(cell)}")
print(f"original_fn:      {original_fn}")
print(f"type(original):   {type(original_fn)}")
print()
print(f"original name:    {original_fn.__name__!r}")
print(f"original doc:     {original_fn.__doc__!r}")
print(f"original result:  {original_fn('doh')}")
\end{verbatim}

Possible output:

\begin{verbatim}
cell:             <cell at 0x...: function object at 0x...>
type(cell):       <class 'cell'>
original_fn:      <function shout at 0x...>
type(original):   <class 'function'>

original name:    'shout'
original doc:     'Return uppercase text.'
original result:  DOH
\end{verbatim}

The original function was not destroyed. It is just no longer directly bound to the public name \texttt{shout}. The wrapper closure still points to it.

The wrapper's code object records the captured free variable.

\begin{verbatim}
def make_logged(fn):
    def wrapper(text):
        print(f"calling {fn.__name__}")
        return fn(text)

    return wrapper

@make_logged
def shout(text):
    """Return uppercase text."""
    return text.upper()

print(f"freevars: {shout.__code__.co_freevars}")

for name, cell in zip(shout.__code__.co_freevars, shout.__closure__):
    print(f"{name:<8} -> {cell.cell_contents}")
\end{verbatim}

Possible output:

\begin{verbatim}
freevars: ('fn',)
fn       -> <function shout at 0x...>
\end{verbatim}

The wrapper function has a free variable named \texttt{fn}. Its closure cell contains the original function object.

A wrapper should usually accept flexible arguments so it can forward the same call shape to the original function.

A too-specific wrapper is fragile.

\begin{verbatim}
def bad_logged(fn):
    def wrapper(text):
        print(f"calling {fn.__name__}")
        return fn(text)

    return wrapper

@bad_logged
def build_line(name, score=42):
    return f"{name:<10} score={score:04d}"

print(build_line("Homer"))

try:
    print(build_line("Lisa", 100))
except TypeError as err:
    print(f"error type: {type(err)}")
    print(f"message:    {err}")
\end{verbatim}

Possible output:

\begin{verbatim}
calling build_line
Homer      score=0042
error type: <class 'TypeError'>
message:    bad_logged.<locals>.wrapper() takes 1 positional argument but 2 were given
\end{verbatim}

The original function could accept two arguments, but the wrapper could accept only one.

A flexible wrapper uses \texttt{*args} and \texttt{**kwargs}.

\begin{verbatim}
def logged(fn):
    def wrapper(*args, **kwargs):
        print(f"calling {fn.__name__}")
        print(f"args:   {args}")
        print(f"kwargs: {kwargs}")

        result = fn(*args, **kwargs)

        print(f"result: {result!r}")
        return result

    return wrapper

@logged
def build_line(name, score=42, *, mood="calm"):
    return f"{name:<10} score={score:04d} mood={mood}"

build_line("Homer")
print()
build_line("Lisa", 100, mood="focused")
\end{verbatim}

Possible output:

\begin{verbatim}
calling build_line
args:   ('Homer',)
kwargs: {}
result: 'Homer      score=0042 mood=calm'

calling build_line
args:   ('Lisa', 100)
kwargs: {'mood': 'focused'}
result: 'Lisa       score=0100 mood=focused'
\end{verbatim}

The wrapper receives arbitrary positional and keyword arguments, then forwards them to the original function with:

\begin{verbatim}
fn(*args, **kwargs)
\end{verbatim}

This preserves the call boundary shape.

The \texttt{args} object is a tuple and \texttt{kwargs} is a dictionary.

\begin{verbatim}
def inspect_call(fn):
    def wrapper(*args, **kwargs):
        print(f"type(args):   {type(args)}")
        print(f"type(kwargs): {type(kwargs)}")

        return fn(*args, **kwargs)

    return wrapper

@inspect_call
def say(text, *, times=1):
    return text * times

print(say("Ni!", times=2))
\end{verbatim}

Possible output:

\begin{verbatim}
type(args):   <class 'tuple'>
type(kwargs): <class 'dict'>
Ni!Ni!
\end{verbatim}

The wrapper is still a normal function with a variadic signature.

The remaining issue is metadata. After decoration, the public function name points to the wrapper, so metadata such as \texttt{\_\_name\_\_}, \texttt{\_\_doc\_\_}, and annotations may describe the wrapper rather than the original function.

The standard solution is \texttt{functools.wraps}.

\begin{verbatim}
import functools

def logged(fn):
    @functools.wraps(fn)
    def wrapper(*args, **kwargs):
        print(f"calling {fn.__name__}")
        return fn(*args, **kwargs)

    return wrapper

@logged
def build_line(name: str, score: int = 42) -> str:
    """Build one formatted display line."""
    return f"{name:<10} score={score:04d}"

print(build_line("Homer"))
print()
print(f"build_line:                 {build_line}")
print(f"build_line.__name__:        {build_line.__name__!r}")
print(f"build_line.__doc__:         {build_line.__doc__!r}")
print(f"build_line.__annotations__: {build_line.__annotations__}")
print(f"build_line.__wrapped__:     {build_line.__wrapped__}")
\end{verbatim}

Possible output:

\begin{verbatim}
calling build_line
Homer      score=0042

build_line:                 <function build_line at 0x...>
build_line.__name__:        'build_line'
build_line.__doc__:         'Build one formatted display line.'
build_line.__annotations__: {'name': <class 'str'>, 'score': <class 'int'>, 'return': <class 'str'>}
build_line.__wrapped__:     <function build_line at 0x...>
\end{verbatim}

The wrapper is still the object bound to \texttt{build\_line}, but \texttt{functools.wraps} copies useful metadata from the original function onto the wrapper. It also stores the original function under \texttt{\_\_wrapped\_\_}.

The \texttt{functools.wraps} object can be inspected.

\begin{verbatim}
import functools

print(f"functools.wraps:       {functools.wraps}")
print(f"type(functools.wraps): {type(functools.wraps)}")
print()

selected_names = [
    "__call__",
    "__name__",
    "__doc__",
    "__module__",
]

for name in selected_names:
    print(f"{name:<14} {name in dir(functools.wraps)}")
\end{verbatim}

Possible output:

\begin{verbatim}
functools.wraps:       <function wraps at 0x...>
type(functools.wraps): <class 'function'>

__call__      True
__name__      True
__doc__       True
__module__    True
\end{verbatim}

\texttt{functools.wraps(fn)} is itself a decorator factory. It returns a decorator that modifies the wrapper function's metadata.

The structure is:

\begin{verbatim}
@functools.wraps(fn)
def wrapper(...):
    ...
\end{verbatim}

which means:

\begin{verbatim}
def wrapper(...):
    ...

wrapper = functools.wraps(fn)(wrapper)
\end{verbatim}

So even \texttt{wraps} follows the same decorator rule: create function object, pass it through a decorator, rebind the name to the returned object.

A minimal comparison:

\begin{verbatim}
import functools

def logged_without_wraps(fn):
    def wrapper(*args, **kwargs):
        return fn(*args, **kwargs)

    return wrapper

def logged_with_wraps(fn):
    @functools.wraps(fn)
    def wrapper(*args, **kwargs):
        return fn(*args, **kwargs)

    return wrapper

@logged_without_wraps
def first():
    """First docstring."""
    return "D'oh!"

@logged_with_wraps
def second():
    """Second docstring."""
    return "No soup for you!"

print(f"first.__name__:  {first.__name__!r}")
print(f"first.__doc__:   {first.__doc__!r}")
print()
print(f"second.__name__: {second.__name__!r}")
print(f"second.__doc__:  {second.__doc__!r}")
print(f"has __wrapped__: {'__wrapped__' in dir(second)}")
\end{verbatim}

Possible output:

\begin{verbatim}
first.__name__:  'wrapper'
first.__doc__:   None

second.__name__: 'second'
second.__doc__:  'Second docstring.'
has __wrapped__: True
\end{verbatim}

The wrapper still controls execution, but the metadata now reflects the original function more accurately.

This matters for documentation, debugging, stack traces, introspection, help output, decorators stacked on decorators, and tools that inspect function signatures.

The final decorator model is:

\begin{verbatim}
def decorator(fn):
    @functools.wraps(fn)
    def wrapper(*args, **kwargs):
        before behavior
        result = fn(*args, **kwargs)
        after behavior
        return result

    return wrapper

@decorator
def target(...):
    original body
\end{verbatim}

and the runtime meaning is:

\begin{verbatim}
target = decorator(target)
\end{verbatim}

The final rule is:

\begin{quote}
A decorator is a definition-time function transformation. The \texttt{@} syntax performs function-object passing and name rebinding. A wrapper usually preserves access to the original function through a closure, and \texttt{functools.wraps} preserves the original function's public metadata on the wrapper.
\end{quote}